\section{Agent Swarm 与协同模式}

\subsection{Swarm 结构分层}

Omar 提到的 Swarm 通常包含 three tiers：1) CLI/Coordinator, 2) execution agents, 3) verification agents。Coordinator 负责 orchestration，execution agents 处理 prompt/LLM calls，verification agents 检查 outputs 并给出 `evidence rating`。

\begin{importantbox}{层级分离的价值}
层级分离降低了复杂度。即使某个 execution agent 输出失误，verification agent 也可以在 `Evidence Matrix` 级别阻断其通过 pipeline。
\end{importantbox}

\subsection{工具链协作}

Swarm 中的 tools 通常以 `MCP` 形式暴露，DSPy 在定义 Module 时会注册 `ToolDescriptor`。在 runtime，Coordinator 会根据 `signature` 顺序调用 tools，并在 evidence log 中记录调用参数与返回值。

\begin{itemize}
    \item `search\_tool`：快速调用 retrieval index
    \item `cohere\_tool`：调用 LLM for consistent reasoning
    \item `validate\_tool`：对 tool output run schema + rule-based checks
\end{itemize}

\subsection{本章小结}

Agent Swarm 的成功依赖于层级结构与工具协调，Coordinator/Validator/Executor 各司其职才能把 compound system 控制住。

